%!TeX root=00_Abschlussarbeit.tex
% \pdfcompresslevel=0
% \pdfobjcompresslevel=0

% Pakete für die Darstellung und Eingabemöglichkeit von Umlauten
\usepackage[T1]{fontenc}
\usepackage[utf8]{inputenc}
\usepackage{lmodern}
% \usepackage{fontspec}   ← auskommentiert lassen
% \setmainfont{...}       ← auskommentiert lassen


%Für Tabellen
\usepackage{rotating}   % Drehen
\usepackage{colortbl}   % Farben
\usepackage{booktabs}   % dickere Linien
\usepackage{longtable}  % für lange Tabellen
\usepackage{tabularx}
\usepackage{array}
\usepackage{multirow}
\usepackage{diagbox}    % für diagonale Linien in Tabellen
\usepackage{float}      % für Tabellenpositionierung [H]
\usepackage{xcolor, colortbl}
\usepackage{caption}    % für \captionof
\usepackage{array}
\newcolumntype{L}[1]{>{\raggedright\arraybackslash}p{#1}}
\newcolumntype{C}[1]{>{\centering\arraybackslash}p{#1}}




% Benötigte Pakete in der Präambel:
 \usepackage{booktabs}
 \usepackage{array}
 \usepackage{colortbl}
 \usepackage[table]{xcolor}
 \usepackage{graphicx}
 \usepackage{diagbox}
 \usepackage{threeparttable}


% Anpassung Seitenränder
\usepackage[left= 3cm, right = 2cm, bottom = 3 cm, top = 2cm]{geometry} 

% Verwendung deutscher Begriffe für automatisch generierte Worte wie z.B. "Inhaltsverzeichnis"
\usepackage[ngerman, english]{babel}

% Anführungszeichen/Zitate
\usepackage{csquotes}

% Trennen wenn Latex es nicht richtig macht:
% \hyphenation{}

% Paket um Dummy-Text/Blindtext zu erzeugen
\usepackage{blindtext}

% Eigenes Titelblatt
\usepackage{LTXKursTitel}

% Zeilenabstand auf 1,5 setzen
\usepackage[onehalfspacing]{setspace}
    \AfterTOCHead{\singlespacing}
    %\KOMAoptions{DIV=last}          % KOMA Klasse für europ. Layout
\usepackage{scrdate, scrtime}       % Zeit und Datumsbefehle
\usepackage{scrlayer-scrpage}       % Erweitere Layout-Optionen
    \pagestyle{scrheadings}         % Seitenlayout selbst definieren

% Belegung von KOMA Variablen für Kopf- und Fußzeilen Gestaltung
    \newcommand{\footlinetext}{\footnotesize \textsf{\color{gray} Kapitel \thesection \  \normalsize}}
    \KOMAoptions{headsepline = yes, footsepline = no}
    \ihead{}
    \chead{Kapitel \thesection}
    \ohead{}
    \ifoot{}
    \cfoot{\pagemark}
    \ofoot{}

% Pakete zum Einbinden von Bildern und Farben
\usepackage{xcolor}
\definecolor{chapterblue}{RGB}{0,70,140}
\setkomafont{chapter}{\color{chapterblue}\huge\sffamily\bfseries}
\setkomafont{section}{\color{chapterblue}\Large\sffamily\bfseries}
\setkomafont{subsection}{\color{chapterblue}\large\sffamily\bfseries}
\setkomafont{subsubsection}{\color{chapterblue}\large\sffamily\bfseries}

 

\usepackage{graphicx}
\usepackage{pdfpages} %PDFs einbinden


% Pakete für Mathematik
\usepackage[free-standing-units,locale = DE]{siunitx}   % Befehle für SI-Einheiten
\usepackage{amsmath}                                    % Mathematik Befehle
\usepackage{amsfonts}

% Paket und Spezifikation der Parameter zur Darstellung von Programmcode mit Courier-Schriftart
\usepackage{listings}  % Darstellung von Quellcode
\usepackage{courier}   % Schriftart laden
    \lstset{
    language=Python,                % Programmiersprache
    basicstyle=\footnotesize\ttfamily,  % Standardschrift
    numbers=left,                       % Ort der Zeilennummern
    numberstyle=\tiny,                  % Stil der Zeilennummern
    %stepnumber=2,  % Abstand zwischen den Zeilennummern
    numbersep=5pt,  % Abstand der Nummern zum Text
    tabsize=2,      % Groesse von Tabs
    extendedchars=true, %
    breaklines=true,    % Zeilen werden Umgebrochen
    keywordstyle=\color{blue}\bfseries,
    frame=b,
    % keywordstyle=[1]\textbf, % Stil der Keywords
    % keywordstyle=[2]\textbf, %
    % keywordstyle=[3]\textbf, %
    % keywordstyle=[4]\textbf, \sqrt{\sqrt{}} %
    stringstyle=\color{magenta}\ttfamily, % Farbe der String
    showspaces=false,   % Leerzeichen anzeigen ?
    showtabs=false,     % Tabs anzeigen ?
    %xleftmargin=17pt,  % Abstände
    %framexleftmargin=17pt,
    %framexrightmargin=5pt,
    %framexbottommargin=4pt,
    commentstyle=\color{gray},
    %backgroundcolor=\color{grey},
    showstringspaces=false,      % Leerzeichen in Strings anzeigen ?
    morekeywords={__global__},  % additional language specific keywords
    morecomment=[l]
    }
    \lstloadlanguages{% Check Dokumentation for further languages ...
    %[Visual]Basic
    %Pascal
    C,
    C++,
    %XML
    %HTML
    %Matlab
    %Java
    }

% Pfade und Topics umbrechbar darstellen
\usepackage{url}
\newcommand{\topic}[1]{\path{#1}}


% Literaturverzeichnis mit BibLaTeX und Biber erstellen
\usepackage[style=ieee, sorting=none, backend=biber]{biblatex}
% bib-Datei einbinden
\addbibresource{Literatur.bib}

\usepackage{abstract}

\usepackage[printonlyused]{acronym} %Abkürzungsverzeichnis

\usepackage{tocbasic}
\usepackage{textcmds}
\usepackage{scrhack}    %gets rid of \float@addtolists error

\allowdisplaybreaks[2]


\usepackage{enumitem}
\setlist{itemsep=0.05\baselineskip, leftmargin=1.4cm}

% Punkte im Inhaltsverzeichnis einfügen (mit tocbasic)

% Punkte ins Inhaltsverzeichnis
\DeclareTOCStyleEntry[beforeskip=5pt, entryformat=\normalfont, pagenumberformat=\normalsize, linefill=\dotfill]{chapter}{chapter}
\DeclareTOCStyleEntry[beforeskip=5pt, indent=1.5em, pagenumberformat=\normalsize, linefill=\dotfill]{tocline}{section}
\DeclareTOCStyleEntry[beforeskip=5pt, indent=3.8em, pagenumberformat=\normalsize, linefill=\dotfill]{tocline}{subsection}




\usepackage{hyperref}   % Layout im PDF Viewer. Muss als letztes Paket geladen werden
    \hypersetup{%
    plainpages=false,
    linktocpage=false,%true,
    breaklinks=true,
    colorlinks=true,
    linkcolor=black,%blue,
    anchorcolor=black,
    citecolor=black,%green,
    filecolor=black,%blue,
    urlcolor=black,%blue%
    pdfstartview={FitV},
    pdfview={FitH},
    pdfpagelayout={SinglePage},
    %pdfpagemode={None},   %unknown value "None"
    }

\usepackage[xindy, order=letter, acronym=true, toc]{glossaries} % muss nach hyperref geladen werden
% \usepackage{glossaries-german} % verwendete Sprachen laden, https://www.ctan.org/pkg/glossaries-german, nötig?

% Zusätzliches Symbolverzeichnis definieren, wichtig vor \makeglossaries
\newglossary[slg]{symbols}{syi}{syg}{Symbolverzeichnis}


\makeglossaries

% Stil definieren (dein eigener Stil)
% Wichtig: description-Umgebung braucht \item, sonst kommt "missing \item".
\usepackage{longtable}

% Zuordnung der Buchstaben-Gruppen (aus den sort-Schlüsseln in 04_Symbole.tex)
% zu den Kategorieüberschriften des Symbolverzeichnisses.
% Unbekannte Buchstaben werden unverändert ausgegeben (Fallback).
\makeatletter
\newcommand*{\symgroupname}[1]{%
  \ifcsname symgrp@#1\endcsname\csname symgrp@#1\endcsname\else#1\fi
}
\@namedef{symgrp@A}{Typografische Konvention}
\@namedef{symgrp@B}{Koordinatensysteme und Frames}
\@namedef{symgrp@C}{Orientierungsdarstellung}
\@namedef{symgrp@D}{Koordinatentransformationen}
\@namedef{symgrp@E}{Extrinsische Kalibrierung}
\@namedef{symgrp@F}{Kamera-Kalibrierung}
\@namedef{symgrp@G}{Extended Kalman Filter}
\@namedef{symgrp@H}{Faktorgraph-Optimierung}
\@namedef{symgrp@I}{Posenberechnung je Marker}
\@namedef{symgrp@J}{Kartenunsicherheit in der Gewichtung}
\@namedef{symgrp@K}{Ausreißerfilterung mit RANSAC}
\@namedef{symgrp@L}{Kamerafusion und Glättung}
\@namedef{symgrp@M}{Kovarianzschätzung}
\@namedef{symgrp@N}{Dead-Reckoning-Lokalisierung}
\makeatother

% Kompakter, tabellenartiger Stil für das Symbolverzeichnis
\newglossarystyle{symstyle}{
  \renewenvironment{theglossary}
    {\begin{longtable}{@{}p{1.3cm}p{\dimexpr\linewidth-1.6cm\relax}@{}}} % Abstand zwischen Symbol und Beschreibung
    {\end{longtable}}

  \renewcommand*{\glossentry}[2]{%
    \glstarget{##1}{\glossentryname{##1}} &
    \glossentrydesc{##1} \\[0.3em] % Abstand zwischen Einträgen
  }

  \renewcommand*{\subglossentry}[3]{%
    \glossentry{##2}{##3}%
  }

  % Benannte Kategorieüberschriften (Titel via \symgroupname aus dem Gruppen-Buchstaben)
  \renewcommand*{\glsgroupheading}[1]{%
    \multicolumn{2}{@{}l}{%
      \rule{0pt}{4.2ex}%
      \normalfont\sffamily\bfseries\large\color{chapterblue}\symgroupname{##1}%
    }\\[0.5em]%
  }
  \renewcommand*{\glsgroupskip}{}
}


% Eigener Stil für das Abkürzungsverzeichnis
\newglossarystyle{acronymstyle}{
  \renewenvironment{theglossary}
    {\begin{longtable}{@{}p{3cm}p{\dimexpr\linewidth-3cm\relax}@{}}}
    {\end{longtable}}

  \renewcommand*{\glossentry}[2]{%
    \glstarget{##1}{\glossentryname{##1}} &
    \glossentrydesc{##1} \\[0.3em] % Abstand zwischen Einträgen
  }

  \renewcommand*{\subglossentry}[3]{%
    \glossentry{##2}{##3}%
  }

  \renewcommand*{\glsgroupheading}[1]{}
  \renewcommand*{\glsgroupskip}{}
}


\setacronymstyle{long-short}
\setglossarystyle{list} %alternativ: tree, list, long3col, index weitere: https://www.dickimaw-books.com/gallery/glossaries-styles/#long

\setcounter{tocdepth}{4} %Erweiterung des Inhaltsverzeichnisses um eine weitere Ebene
\setcounter{secnumdepth}{4} %Erweiterung der Kapiteltiefe um eine weitere Ebene


\renewcommand{\familydefault}{\sfdefault}
%\usepackage{showframe}

% Kopfzeilen-Abstand reduzieren
\setlength{\headheight}{14pt}
\setlength{\headsep}{10pt}
 
% Kapitel oben beginnen lassen
\RedeclareSectionCommand[
  beforeskip=0pt,
  afterskip=1\baselineskip
]{chapter}
 
\RedeclareSectionCommand[
  beforeskip=1.5ex plus 0.5ex minus 0.2ex,
  afterskip=0.5ex
]{section}
 
\RedeclareSectionCommand[
  beforeskip=1ex plus 0.5ex minus 0.2ex,
  afterskip=0.5ex
]{subsection}

 
 
% Abstract nicht mittig
\setlength{\absleftindent}{0pt}
\setlength{\absrightindent}{0pt}
\setlength{\abstitleskip}{0pt}


% Einheitliche Abstände um abgesetzte Gleichungen
\AtBeginDocument{%
  \setlength{\abovedisplayskip}{0.7\baselineskip}
  \setlength{\belowdisplayskip}{0.7\baselineskip}
  \setlength{\abovedisplayshortskip}{0.7\baselineskip}
  \setlength{\belowdisplayshortskip}{0.7\baselineskip}
}


%%%%%%%%%%%%%%%%%%%%%%%%%%%%%%%%%%%%%%%%%%%%%%%%%%%
%%%          Variablen zum Belegen              %%%
%%%%%%%%%%%%%%%%%%%%%%%%%%%%%%%%%%%%%%%%%%%%%%%%%%%

% Variables for title page
    \DAAutor{\par Hannes Müller\par 3796115}
    \DATyp{}
    \DAAutorAdresse{Masterarbeit im Studiengang Elektronische und Mechatronische Systeme\par Vertiefungsrichtung Automatisierungstechnik}
    \DAFachbereich{Elektrotechnik Feinwerktechnik Informationstechnik}
    \DATitel{Implementierung eines fiducialbasierten Lokalisierungssystems für Luftfahrzeuge im Innenraum}
    \DABetreuerTextA{Betreuer}
    \DABetreuerTextB{Betreuer}
    \DABetreuerA{Prof. Dr. rer. nat. Pfitzner}
    \DABetreuerB{Prof. Dr.-Ing. Schröder}
    \DAOrt{Abgabe: Nürnberg}
    \DAAbgabedatum{??.\,Juli 2026}
    \DAAbgabesemester{Sommersemester 2026}